\subsection{Modes of Convergence}

% **Role:** compare pointwise, almost-everywhere, uniform, almost-uniform, in-measure, and $L^p$ convergence.

Throughout this subsection $(X,\mathcal{M},\mu)$ is a fixed measure space and, unless otherwise stated, functions are real-valued and measurable. The point is not to rank convergence modes linearly, but to know which hypotheses convert one mode into another.

\begin{definition} \label{def:modes-pointwise-uniform-ae} \hlt{Pointwise, Uniform, and Almost-Everywhere Convergence}\\
Let $(f_n)$ and $f$ be functions on $X$.
\begin{enumerate}[label=(\roman*)]
\setlength{\itemsep}{0pt}
\item $f_n\to f$ \hlt{pointwise} if $\forall x\in X$, $\varepsilon>0$, $\exists N(\varepsilon,x)$ s.t. $n\geq N(\varepsilon,x)$ implies $\abs{f_n(x)-f(x)}<\varepsilon$.
\item $f_n\to f$ \hlt{uniformly} if $\forall \varepsilon>0$ $\exists N(\varepsilon)$ s.t. $n\geq N(\varepsilon)$ implies $\abs{f_n(x)-f(x)}<\varepsilon$ $\forall x\in X$.
\item $f_n\to f$ \hlt{almost everywhere} if $\exists A\in\mathcal{M}$ with $\mu(A)=0$ s.t. $\forall x\in X\setminus A$, $\varepsilon>0$,\\
$\exists N(\varepsilon,x)$ s.t. $n\geq N(\varepsilon,x)$ implies $\abs{f_n(x)-f(x)}<\varepsilon$.
\end{enumerate}
\end{definition}

\begin{proposition} \label{prop:basic-mode-implications} \hlt{First Implications and First Failures}\\
Let $f_n,f:X\to\R$.
\begin{enumerate}[label=(\roman*)]
\setlength{\itemsep}{0pt}
\item $f_n\rightrightarrows f\Rightarrow f_n\to f$ pointwise $\Rightarrow f_n\to f$ a.e.; converses fail.
\item $\abs{X}<+\infty$ and $f_n\to f$ pointwise $\Rightarrow f_n\rightrightarrows f$.
\item $\mu(A)=0\Rightarrow A=\emptyset$ and $f_n\to f$ a.e. $\Rightarrow f_n\to f$ pointwise.
\end{enumerate}
\begin{proof}
\begin{enumerate}[label=(\roman*)]
\setlength{\itemsep}{0pt}
\item If $f_n\rightrightarrows f$, then $\forall\varepsilon>0$ $\exists N_\varepsilon$ such that $n\geq N_\varepsilon$ implies $\abs{f_n(x)-f(x)}<\varepsilon$ for all $x\in X$; hence pointwise convergence with $N(\varepsilon,x)=N_\varepsilon$. Pointwise convergence is a.e. convergence with $A=\emptyset$. Conversely, on $([0,1],\lambda)$, $x^n\to\mathbf{1}_{\{1\}}$ pointwise but not uniformly, and $\mathbf{1}_{\{0\}}\to0$ a.e. but not pointwise.
\item If $X$ is finite and $f_n\to f$ pointwise, then for each $\varepsilon>0$ and $x\in X$ choose $N(\varepsilon,x)$.\\
Set $N_\varepsilon=\max_{x\in X}N(\varepsilon,x)$; then $n\geq N_\varepsilon$ implies $\abs{f_n(x)-f(x)}<\varepsilon$ for all $x\in X$, so $f_n\rightrightarrows f$.
\item If $f_n\to f$ a.e., then $\exists A\in\mathcal{M}$ with $\mu(A)=0$ and $f_n(x)\to f(x)$ for all $x\in X\setminus A$.\\
Since every null set is empty, $A=\emptyset$, so $f_n\to f$ pointwise.
\end{enumerate}
\hlt{Structural Commentary}: Uniformity is bounded $x$-dependence of $N(\varepsilon,x)$.\\
\hlt{Key Insight}: Uniformity controls indices; a.e. convergence controls null exceptions.
\end{proof}
\end{proposition}

\begin{proposition} \label{prop:uniform-to-lp-on-finite-measure} \hlt{Uniform Convergence Gives $L^p$ Convergence on Finite Measure Spaces}\\
Let $1\leq p<\infty$, $\mu(X)<+\infty$, and $f_n\in L^p(\mu)$. If $f_n\rightrightarrows f$ on $X$, then $f\in L^p(\mu)$ and $\norm{f_n-f}_p\to0$.
\begin{proof}
For $\varepsilon>0$, choose $N$ with $n\geq N\Rightarrow \abs{f_n-f}\leq\varepsilon$ on $X$.\\
Then $n\geq N\Rightarrow \norm{f_n-f}_p\leq(\int_X\varepsilon^p\,d\mu)^{1/p}=\varepsilon\mu(X)^{1/p}$, so $\norm{f_n-f}_p\to0$.\\
Taking $\varepsilon=1$, choose $n_0$ with $\abs{f_{n_0}-f}\leq1$ on $X$; hence $f-f_{n_0}\in L^p(\mu)$ and $f=f_{n_0}+(f-f_{n_0})\in L^p(\mu)$.\\
\hlt{Structural Commentary}: Uniform error becomes $L^p$ error after multiplication by $\mu(X)^{1/p}$.\\
\hlt{Key Insight}: Finite measure converts a pointwise sup-bound into an integral bound.
\end{proof}
\end{proposition}

\begin{theorem} \label{thm:dominated-ae-to-lp} \hlt{Dominated Almost-Everywhere Convergence Gives $L^p$ Convergence}\\
Let $1\leq p<\infty$, $f_n\in L^p(\mu)$, and $f_n\to f$ a.e.\\
If $\exists g\in L^p(\mu)$ such that $\abs{f_n}\leq g$ a.e. for all $n$, then $f\in L^p(\mu)$ and $\norm{f_n-f}_p\to0$.
\begin{proof}
Off a null set, $f_n\to f$ and $\abs{f_n}\leq g$ for all $n$; hence $\abs f\leq g$, so $f\in L^p(\mu)$. Also $\abs{f_n-f}^p\to0$ a.e. and $\abs{f_n-f}^p\leq(2g)^p\in\mathscr{L}^1(\mu)$ a.e.; Theorem \ref{thm:dominated-convergence-lebesgue} gives $\norm{f_n-f}_p^p=\int_X\abs{f_n-f}^p\,d\mu\to0$.\\
\hlt{Structural Commentary}: Apply ordinary dominated convergence to the $p$th-power error.\\
\hlt{Key Insight}: A single $L^p$ envelope controls both membership of the limit and convergence in mean.
\end{proof}
\end{theorem}

\begin{corollary} \label{cor:bounded-ae-finite-measure-lp} \hlt{Uniform Boundedness on Finite Measure Spaces}\\
Let $\mu(X)<+\infty$, $0\leq K<+\infty$, and $f_n\to f$ a.e.\\
If $\abs{f_n}\leq K$ a.e. for all $n$, then for every $1\leq p<\infty$, $f\in L^p(\mu)$ and $\norm{f_n-f}_p\to0$.
\begin{proof}
Fix $1\leq p<\infty$ and set $g\defeq K\mathbf{1}_X$. Then $\int_Xg^p\,d\mu=K^p\mu(X)<+\infty$, so $g\in L^p(\mu)$, and $\abs{f_n}\leq g$ a.e. for all $n$. Since $f_n\to f$ a.e., Theorem \ref{thm:dominated-ae-to-lp} gives $f\in L^p(\mu)$ and $\norm{f_n-f}_p\to0$.\\
\hlt{Structural Commentary}: Finite measure turns a uniform a.e. bound into an $L^p$ envelope.\\
\hlt{Key Insight}: Bounded a.e. convergence on finite measure is dominated convergence.
\end{proof}
\end{corollary}

\begin{example} \label{ex:modes-uniform-not-lp-infinite-measure} \hlt{Uniform Need Not Imply $L^p$ on Infinite Measure}\\
On $(\R,\mathcal{B}(\R),\lambda)$, fix $1\leq p<\infty$ and set $f_n\defeq n^{-1/p}\mathbf{1}_{[0,n]}$. Then $\norm{f_n}_\infty=n^{-1/p}\to0$, so $f_n\rightrightarrows0$.\\
However $\norm{f_n}_p^p=\int_0^n n^{-1}\,d\lambda=1$, hence $\norm{f_n}_p\nrightarrow0$.\\
\hlt{Structural Commentary}: The support grows exactly enough to preserve $L^p$ mass.\\
\hlt{Key Insight}: Uniform smallness alone does not control integrals on infinite measure spaces.
\end{example}

\begin{example} \label{ex:modes-ae-measure-not-lp-finite-measure} \hlt{Almost Everywhere and In-Measure Need Not Imply $L^p$}\\
On $[0,1]$ with Lebesgue measure, set $f_n\defeq n\mathbf{1}_{[1/n,\,2/n]}$.\\
For each $x\in[0,1]$, $x\in[1/n,2/n]$ for only finitely many $n$, hence $f_n(x)\to0$.\\
If $\alpha>0$ and $n\geq\alpha$, then $\lambda\{\abs{f_n}\geq\alpha\}=1/n\to0$, so $f_n\to0$ in measure.\\
Yet for $1\leq p<\infty$, $\norm{f_n}_p^p=n^p\lambda([1/n,2/n])=n^{p-1}\nrightarrow0$.\\
Given $\delta>0$, choose $0<\eta<\delta$; on $[0,1]\setminus[0,\eta]$, $f_n=0$ for all $n>2/\eta$, so $f_n\to0$ almost uniformly.\\
Thus a.e., in-measure, and almost-uniform convergence do not imply $L^p$ convergence.
\end{example}

\begin{figure}[H]
  \centering
  % Source: attached Chapter 7, Example 7.4; clean redraw of the interval
% enumeration behind the typewriter sequence.
\begin{tikzpicture}[
  x=8.6cm,
  y=0.78cm,
  >=Stealth,
  every node/.style={font=\scriptsize},
  base/.style={black!65, line width=0.55pt},
  tick/.style={black!45, line width=0.45pt},
  interval/.style={fill=blue!8, draw=blue!45!black, line width=0.5pt},
  active/.style={fill=blue!28, draw=blue!70!black, line width=0.7pt},
  hit/.style={red!70!black, dashed, line width=0.7pt}
]
  \def\xhit{0.63}
  \foreach \m/\y in {1/4,2/3,3/2,4/1,5/0} {
    \draw[base] (0,\y) -- (1,\y);
    \node[left=7pt] at (0,\y) {$m=\m$};
    \pgfmathtruncatemacro{\hitj}{floor(\m*\xhit)+1}
    \foreach \j in {1,...,\m} {
      \pgfmathsetmacro{\a}{(\j-1)/\m}
      \pgfmathsetmacro{\b}{\j/\m}
      \ifnum\j=\hitj
        \filldraw[active] (\a,\y-0.13) rectangle (\b,\y+0.13);
      \else
        \filldraw[interval] (\a,\y-0.13) rectangle (\b,\y+0.13);
      \fi
      \draw[tick] (\a,\y-0.22) -- (\a,\y+0.22);
    }
    \draw[tick] (1,\y-0.22) -- (1,\y+0.22);
  }

  \draw[hit] (\xhit,-0.35) -- (\xhit,4.38);
  \node[red!70!black, above] at (\xhit,4.38) {fixed $x$};
  \node[below] at (0,-0.34) {$0$};
  \node[below] at (1,-0.34) {$1$};
\end{tikzpicture}

  \caption{Typewriter intervals shrink, but each fixed $x$ is hit once per partition.}
  \label{fig:typewriter-intervals-lp-no-pointwise}
\end{figure}

\begin{example} \label{ex:typewriter-lp-no-pointwise} \hlt{$L^p$ Convergence Without Pointwise Convergence}\\
Let $X=[0,1]$ with Lebesgue measure and list $I_{m,k}\defeq[(k-1)/m,k/m]$, $1\leq k\leq m$, row by row in $m$.\\
Let $f_n\defeq\mathbf{1}_{I_n}$, $I_n$ is $n$th interval. If $I_n$ in row $m(n)$, then $m(n)\to\infty$, $\norm{f_n}_p^p=\lambda(I_n)=1/m(n)\to0$; hence $f_n\to0$ in $L^p$, and therefore in measure by Proposition \ref{prop:lp-convergence-implies-in-measure}.\\
$\forall x\in[0,1]$, every row contains some $I_{m,k}$ with $x\in I_{m,k}$, giving a subsequence with value $1$ at $x$; for infinitely many rows, some $I_{m,k}$ omits $x$, giving a subsequence with value $0$ at $x$. Thus $(f_n(x))$ converges nowhere.\\
The subfamily $I_{m,2}=[1/m,2/m]$, $m\geq2$, gives a subsequence converging pointwise to $0$ on all of $X$.
\end{example}

\begin{definition} \label{def:convergence-in-measure-and-cauchy} \hlt{Convergence and Cauchy Convergence in Measure}\\
Let $(f_n)$ be measurable functions on $X$.
\begin{enumerate}[label=(\roman*)]
\setlength{\itemsep}{0pt}
\item $f_n\to f$ \hlt{in measure} if $\forall\alpha>0$, $\mu\{x\in X:\abs{f_n(x)-f(x)}\geq\alpha\}\to0$.
\item $(f_n)$ is \hlt{Cauchy in measure} if $\forall\alpha>0$, $\mu\{x\in X:\abs{f_m(x)-f_n(x)}\geq\alpha\}\to0$ as $m,n\to\infty$.
\end{enumerate}
The same definitions result with $>$ in place of $\geq$.
\end{definition}

\begin{proposition} \label{prop:lp-convergence-implies-in-measure} \hlt{$L^p$ Convergence Implies Convergence in Measure}\\
Let $1\leq p<\infty$. If $\norm{f_n-f}_p\to0$, then $f_n\to f$ in measure.
\begin{proof}
For $\alpha>0$, set $E_n(\alpha)\defeq\{\abs{f_n-f}\geq\alpha\}$.\\
Since $\norm{f_n-f}_p^p\geq\int_{E_n(\alpha)}\abs{f_n-f}^p\,d\mu\geq\alpha^p\mu(E_n(\alpha))$, $\mu(E_n(\alpha))\leq\alpha^{-p}\norm{f_n-f}_p^p\to0$.\\
\hlt{Structural Commentary}: Large-error sets are controlled by the integral of the $p$th-power error.\\
\hlt{Key Insight}: This is Chebyshev's inequality in deterministic measure form.
\end{proof}
\end{proposition}

\begin{theorem} \label{thm:riesz-subsequence-cauchy-measure} \hlt{Riesz Subsequence Principle for Convergence in Measure}\\
If $(f_n)$ is Cauchy in measure, then there exist a subsequence $(g_k)$ and a measurable $f$ such that:
\begin{enumerate}[label=(\roman*)]
\setlength{\itemsep}{0pt}
\item $g_k\to f$ a.e.
\item $g_k\to f$ in measure.
\item For every $\varepsilon>0$, there is $E\in\mathcal{M}$ with $\mu(E)<\varepsilon$ such that $g_k\to f$ uniformly on $X\setminus E$.
\end{enumerate}
\begin{proof}
\begin{enumerate}[label=(\roman*)]
\setlength{\itemsep}{0pt}
\item Choose $n_1<n_2<\cdots$ so that $m,n\geq n_k\Rightarrow\mu\{\abs{f_m-f_n}\geq2^{-k}\}<2^{-k}$, and set $g_k\defeq f_{n_k}$.\\
For $E_k\defeq\{\abs{g_{k+1}-g_k}\geq2^{-k}\}$, $\mu(E_k)<2^{-k}$. Set $F_k\defeq\bigcup_{j=k}^{\infty}E_j$.\\
Then $F_k\downarrow$, $\mu(F_k)\leq\sum_{j=k}^{\infty}2^{-j}=2^{-(k-1)}$, and for $i>j\geq k$, $x\notin F_k$,\\
$\abs{g_i(x)-g_j(x)}\leq\sum_{r=j}^{i-1}\abs{g_{r+1}(x)-g_r(x)}<2^{-(j-1)}$. Hence $(g_j)$ is uniformly Cauchy on $X\setminus F_k$ $\forall k$.\\
Let $F\defeq\bigcap_kF_k$, so $\mu(F)=0$. Define $f=\lim_jg_j$ on $X\setminus F$, $f=0$ on $F$. $f$ is measurable, $g_j\to f$ a.e.
\item Letting $i\to\infty$ gives $j\geq k$, $x\notin F_k\Rightarrow \abs{f(x)-g_j(x)}\leq2^{-(j-1)}\leq2^{-(k-1)}$. Given $\alpha,\varepsilon>0$, choose $k$ with $2^{-(k-1)}<\alpha$ and $\mu(F_k)<\varepsilon$; then $j\geq k\Rightarrow \{\abs{f-g_j}\geq\alpha\}\subseteq F_k$, so $g_j\to f$ in measure.
\item Given $\varepsilon>0$, choose $k$ with $\mu(F_k)<\varepsilon$.\\
The estimate in (ii) gives $\sup_{x\in X\setminus F_k}\abs{f(x)-g_j(x)}\leq2^{-(j-1)}\to0$, so $g_j\to f$ uniformly on $X\setminus F_k$.
\end{enumerate}
\hlt{Structural Commentary}: Cauchy in measure lets us select summably small exceptional sets.\\
\hlt{Key Insight}: Outside the tail union of exceptional sets, the subsequence is uniformly Cauchy.
\end{proof}
\end{theorem}

\begin{corollary} \label{cor:cauchy-in-measure-completion} \hlt{Completeness for Convergence in Measure}\\
$(f_n)$ is Cauchy in measure iff $\exists$ measurable $f$ with $f_n\to f$ in measure.
\begin{proof}
\begin{enumerate}[label=\roman*.]
\setlength{\itemsep}{0pt}
\item[$(\Rightarrow)$] If $(f_n)$ is Cauchy in measure, Theorem \ref{thm:riesz-subsequence-cauchy-measure} gives $f_{n_k}\to f$ in measure for some measurable $f$.\\
For $\alpha>0$, $\{\abs{f_n-f}\geq\alpha\}\subseteq\{\abs{f_n-f_{n_k}}\geq\alpha/2\}\cup\{\abs{f_{n_k}-f}\geq\alpha/2\}$.\\
Choose $k$ large, then $n$ large, to get $f_n\to f$ in measure.
\item[$(\Leftarrow)$] For $\alpha>0$, $\{\abs{f_m-f_n}\geq\alpha\}\subseteq\{\abs{f_m-f}\geq\alpha/2\}\cup\{\abs{f_n-f}\geq\alpha/2\}$.\\
Hence $\mu\{\abs{f_m-f_n}\geq\alpha\}\to0$ as $m,n\to\infty$.
\end{enumerate}
\hlt{Structural Commentary}: The subsequential in-measure limit attracts the whole Cauchy sequence.\\
\hlt{Key Insight}: Convergence in measure is complete after measurable limits are admitted.
\end{proof}
\end{corollary}

\begin{corollary} \label{cor:cauchy-in-measure-completion-uniqueness} \label{cor:convergence-in-measure-uniqueness} \hlt{Uniqueness for Convergence in Measure}\\
If $f_n\to f$ and $f_n\to g$ in measure, then $f=g$ a.e.
\begin{proof}
For $\alpha>0$, $\{\abs{f-g}\geq\alpha\}\subseteq\{\abs{f-f_n}\geq\alpha/2\}\cup\{\abs{f_n-g}\geq\alpha/2\}$, so $\mu\{\abs{f-g}\geq\alpha\}=0$.\\
Hence $\mu\{\abs{f-g}>0\}=0$ and $f=g$ a.e.\\
\hlt{Structural Commentary}: Two in-measure limits cannot differ on a positive-measure set.\\
\hlt{Key Insight}: In-measure limits are unique modulo null sets.
\end{proof}
\end{corollary}

\begin{proposition} \label{prop:subsequence-and-indicator-exercise-consequences} \hlt{Subsequence and Indicator Consequences}
\begin{enumerate}[label=(\roman*)]
\setlength{\itemsep}{0pt}
\item $f_n\to f$ in measure $\Rightarrow$ every subsequence $f_{n_k}\to f$ in measure; if $(f_n)$ is Cauchy in measure,\\
then every $(f_{n_k})$ is Cauchy in measure.
\item If $\norm{f_n-f}_p\to0$ and $\norm{f_{n_k}-g}_p\to0$, then $f=g$ a.e.
\item If $f_n=\mathbf{1}_{E_n}$ and $\norm{f_n-f}_p\to0$, then $f=\mathbf{1}_E$ a.e. for some $E\in\mathcal{M}$.
\end{enumerate}
\begin{proof}
\begin{enumerate}[label=(\roman*)]
\setlength{\itemsep}{0pt}
\item For $\alpha>0$, $\mu\{\abs{f_{n_k}-f}\geq\alpha\}\to0$ since $n_k\to\infty$; the Cauchy claim is identical with $m_k,n_\ell\to\infty$.
\item Proposition \ref{prop:lp-convergence-implies-in-measure} gives $f_n\to f$ and $f_{n_k}\to g$ in measure; also $f_{n_k}\to f$ in measure by (i).\\
Uniqueness in Corollary \ref{cor:cauchy-in-measure-completion-uniqueness} gives $f=g$ a.e.
\item By Proposition \ref{prop:lp-convergence-implies-in-measure} and Theorem \ref{thm:riesz-subsequence-cauchy-measure}, some $f_{n_k}\to f$ a.e.\\
Since $f_{n_k}\in\{0,1\}$, $f\in\{0,1\}$ a.e.; with $E\defeq\{f=1\}$, $E\in\mathcal{M}$ and $f=\mathbf{1}_E$ a.e.
\end{enumerate}
\hlt{Structural Commentary}: Subsequence permanence and a.e. uniqueness transfer mode limits to representatives.\\
\hlt{Key Insight}: In-measure subsequences preserve limits; indicator limits remain indicators a.e.
\end{proof}
\end{proposition}

\begin{theorem} \label{thm:dominated-measure-to-lp} \hlt{Dominated Convergence in Measure Gives $L^p$ Convergence}\\
Let $1\leq p<\infty$, $f_n\in L^p(\mu)$, $f_n\to f$ in measure, and $\abs{f_n}\leq g$ a.e. for some $g\in L^p(\mu)$.\\
Then $f\in L^p(\mu)$ and $\norm{f_n-f}_p\to0$.
\begin{proof}
Since $f_n\to f$ in measure, $(f_n)$ is Cauchy in measure.\\
By Theorem \ref{thm:riesz-subsequence-cauchy-measure}, some subsequence $(u_j)$ satisfies $u_j\to h$ a.e. and in measure; uniqueness in Corollary \ref{cor:cauchy-in-measure-completion-uniqueness} gives $h=f$ a.e. Since $\abs{u_j}\leq g$ a.e., Theorem \ref{thm:dominated-ae-to-lp} gives $f\in L^p(\mu)$. If $\norm{f_n-f}_p\nrightarrow0$, then $\exists\varepsilon>0$ and a subsequence $(g_k)$ with $\norm{g_k-f}_p\geq\varepsilon$. Since $g_k\to f$ in measure, the same argument gives a subsubsequence $(h_j)$ with $h_j\to f$ a.e.; Theorem \ref{thm:dominated-ae-to-lp} yields $\norm{h_j-f}_p\to0$, contradicting $\norm{h_j-f}_p\geq\varepsilon$.\\
\hlt{Structural Commentary}: Reduce convergence in measure to almost-everywhere convergence along a subsequence.\\
\hlt{Key Insight}: Domination turns every in-measure subsequence into an $L^p$-convergent subsubsequence.
\end{proof}
\end{theorem}

\begin{definition} \label{def:almost-uniform-convergence} \hlt{Almost Uniform Convergence}\\
Let $(f_n)$ be measurable functions on $X$.
\begin{enumerate}[label=(\roman*)]
\setlength{\itemsep}{0pt}
\item $f_n\to f$ \hlt{almost uniformly} if $\forall\delta>0$ $\exists E_\delta\in\mathcal{M}$ with $\mu(E_\delta)<\delta$ and $f_n\rightrightarrows f$ on $X\setminus E_\delta$.
\item $(f_n)$ is \hlt{almost uniformly Cauchy} if $\forall\delta>0$ $\exists E_\delta\in\mathcal{M}$ with $\mu(E_\delta)<\delta$ s.t. $(f_n)$ uniformly Cauchy on $X\setminus E_\delta$.
\end{enumerate}
\end{definition}

\begin{lemma} \label{lem:almost-uniform-cauchy-has-ae-limit} \hlt{Almost Uniform Cauchy Sequences Have Almost Uniform Limits}\\
If $(f_n)$ is almost uniformly Cauchy, then $\exists$ measurable $f$ such that $f_n\to f$ almost uniformly and a.e.
\begin{proof}
Choose $E_k\in\mathcal{M}$ with $\mu(E_k)<2^{-k}$ and $(f_n)$ uniformly Cauchy on $X\setminus E_k$. Set $F_k\defeq\bigcup_{j=k}^{\infty}E_j$ and $F\defeq\bigcap_kF_k$; then $F_k\downarrow$, $\mu(F_k)<2^{-(k-1)}$, $\mu(F)=0$, and $(f_n)$ is uniformly Cauchy on $X\setminus F_k$. By Theorem \ref{thm:uniform-cauchy-criterion}, $f_n$ has a uniform limit on each $X\setminus F_k$; these limits agree on overlaps, so define measurable $f$ on $X\setminus F$, and set $f=0$ on $F$. Hence $f_n\to f$ a.e. and uniformly on each $X\setminus F_k$; since $\mu(F_k)\to0$, $f_n\to f$ almost uniformly.\\
\hlt{Structural Commentary}: Tail unions convert many small exceptional sets into nested uniform domains.\\
\hlt{Key Insight}: Almost uniform Cauchy means uniformly Cauchy after discarding a tail-small exceptional set.
\end{proof}
\end{lemma}

\begin{theorem} \label{thm:almost-uniform-and-measure-subsequences} \hlt{Almost Uniform Convergence and Convergence in Measure}\\
If $f_n\to f$ almost uniformly, then $f_n\to f$ in measure.\\
If $f_n\to f$ in measure, then $\exists(f_{n_k})$ with $f_{n_k}\to f$ almost uniformly.
\begin{proof}
If $f_n\rightrightarrows f$ on $X\setminus E_\delta$ with $\mu(E_\delta)<\delta$, then for each $\alpha>0$ and all large $n$, $\{\abs{f_n-f}\geq\alpha\}\subseteq E_\delta$; hence $\limsup_n\mu\{\abs{f_n-f}\geq\alpha\}\leq\delta$, and $\delta\downarrow0$ gives convergence in measure. Conversely, $f_n\to f$ in measure implies $(f_n)$ is Cauchy in measure; Theorem \ref{thm:riesz-subsequence-cauchy-measure} gives $f_{n_k}\to h$ almost uniformly and in measure. By uniqueness, $h=f$ a.e.; absorbing $\{h\neq f\}$ into the exceptional sets gives $f_{n_k}\to f$ almost uniformly.\\
\hlt{Structural Commentary}: Almost uniform convergence is uniform convergence off arbitrarily small sets.\\
\hlt{Key Insight}: Convergence in measure always hides an almost-uniform subsequence.
\end{proof}
\end{theorem}

\begin{theorem} \label{thm:egoroff-finite-measure} \hlt{Egoroff's Theorem}\\
If $\mu(X)<+\infty$ and $f_n\to f$ a.e., with all functions finite-valued off a null set,\\
then $f_n\to f$ almost uniformly and in measure.
\begin{proof}
Absorb the null exceptional set into every later $E_\delta$, and assume $f_n(x)\to f(x)$ for all $x\in X$.\\
For $m,n\in\N$, set $E_n(m)\defeq\bigcup_{k=n}^{\infty}\{\abs{f_k-f}\geq1/m\}$. Then $E_n(m)\downarrow\emptyset$, so $\mu(E_n(m))\to0$ by Corollary \ref{cor:measure-continuity-up-down}.\\
Given $\delta>0$, choose $k_m$ with $\mu(E_{k_m}(m))<\delta2^{-m}$ and set $E_\delta\defeq\bigcup_{m\geq1}E_{k_m}(m)$.\\
Then $\mu(E_\delta)<\delta$, and $x\notin E_\delta$, $k\geq k_m$ imply $\abs{f_k(x)-f(x)}<1/m$; hence $f_n\rightrightarrows f$ on $X\setminus E_\delta$.\\
The in-measure conclusion follows from Theorem \ref{thm:almost-uniform-and-measure-subsequences}.\\
\hlt{Structural Commentary}: Finite measure turns decreasing bad sets with empty intersection into small bad sets.\\
\hlt{Key Insight}: A.e. convergence is almost uniform on finite measure spaces.
\end{proof}
\end{theorem}

\begin{corollary} \label{cor:finite-measure-ae-implies-measure} \hlt{Finite-Measure Almost-Everywhere Convergence Implies In-Measure Convergence}\\
If $\mu(X)<+\infty$ and $f_n\to f$ a.e., with $f_n,f$ finite-valued off a null set, then $f_n\to f$ in measure.
\begin{proof}
By Theorem \ref{thm:egoroff-finite-measure}, $f_n\to f$ almost uniformly; Theorem \ref{thm:almost-uniform-and-measure-subsequences} gives $f_n\to f$ in measure.\\
\hlt{Structural Commentary}: Finite measure upgrades a.e. convergence through almost uniform convergence.\\
\hlt{Key Insight}: On finite measure spaces, a.e. convergence implies in-measure convergence.
\end{proof}
\end{corollary}

\begin{example} \label{ex:egoroff-needs-finite-valued-limit} \hlt{Egoroff Needs a Finite-Valued Limit}\\
On $([0,1],\lambda)$, set $f_n\defeq n\mathbf{1}_{[0,1]}$. Then $f_n(x)\to+\infty$ for all $x$, so the pointwise limit is not finite-valued. If $f_n\to f$ almost uniformly for finite-valued measurable $f$, then for some $E$ with $\lambda(E)<1/2$, $f_n\rightrightarrows f$ on $[0,1]\setminus E$; this nonempty set has $f_n(x)=n\to+\infty$, contradicting finite-valued pointwise convergence.
\end{example}

% Source: attached Chapter 7, Diagrams 7.1--7.3. Solid arrows are implications;
% dashed arrows mean a subsequence converges in the indicated mode.
\begin{figure}[H]
  \centering
  % Source: attached Chapter 7, Diagrams 7.1--7.3; simplified redraw.
% Solid arrows are implications. Dashed arrows mean a subsequence converges in
% the indicated mode.
\begin{tikzpicture}[
  scale=0.94,
  >=Stealth,
  every node/.style={font=\scriptsize},
  mode/.style={draw=blue!60!black, fill=blue!6, rounded corners=2pt,
    text=black,
    minimum width=0.72cm, minimum height=0.42cm, inner sep=2pt},
  title/.style={font=\scriptsize\itshape, black!75},
  imp/.style={->, black!78, line width=0.8pt},
  eqimp/.style={<->, black!78, line width=0.8pt},
  sub/.style={->, red!65!black, dashed, line width=0.75pt}
]
  \begin{scope}[xshift=0cm]
    \node[title] at (0,1.50) {general};
    \node[mode] (gAE) at (-1.05,0.62) {$AE$};
    \node[mode] (gAU) at (1.05,0.62) {$AU$};
    \node[mode] (gLP) at (-1.05,-0.62) {$L^p$};
    \node[mode] (gM) at (1.05,-0.62) {$M$};

    \draw[imp] (gAU) -- (gAE);
    \draw[imp] (gAU) -- (gM);
    \draw[imp] (gLP) -- (gM);
    \draw[sub] (gM) to[bend left=18] (gAU);
  \end{scope}

  \begin{scope}[xshift=3.65cm]
    \node[title] at (0,1.50) {finite measure};
    \node[mode] (fAE) at (-1.05,0.62) {$AE$};
    \node[mode] (fAU) at (1.05,0.62) {$AU$};
    \node[mode] (fLP) at (-1.05,-0.62) {$L^p$};
    \node[mode] (fM) at (1.05,-0.62) {$M$};

    \draw[eqimp] (fAE) -- (fAU);
    \draw[imp] (fAU) -- (fM);
    \draw[imp] (fLP) -- (fM);
    \draw[sub] (fM) to[bend left=18] (fAU);
  \end{scope}

  \begin{scope}[xshift=7.30cm]
    \node[title] at (0,1.50) {dominated by $g\in L^p$};
    \node[mode] (dAE) at (-1.05,0.62) {$AE$};
    \node[mode] (dAU) at (1.05,0.62) {$AU$};
    \node[mode] (dLP) at (-1.05,-0.62) {$L^p$};
    \node[mode] (dM) at (1.05,-0.62) {$M$};

    \draw[imp] (dAU) -- (dAE);
    \draw[imp] (dAE) -- (dLP);
    \draw[eqimp] (dLP) -- (dM);
    \draw[sub] (dM) to[bend left=18] (dAU);
  \end{scope}
\end{tikzpicture}

  \caption{Modes of convergence. Solid arrows are implications; dashed arrows give a convergent subsequence.}
  \label{fig:modes-convergence-implication-diagrams}
\end{figure}

\begin{method} \label{meth:modes-of-convergence-counterexample-dictionary} \hlt{Counterexample Dictionary}
\begin{enumerate}[label=(\roman*)]
\setlength{\itemsep}{0pt}
\item On $(\R,\lambda)$, $f_n=n^{-1/p}\mathbf{1}_{[0,n]}$ satisfies $f_n\rightrightarrows0$ but $\norm{f_n}_p=1$.
\item On $([0,1],\lambda)$, $f_n=n\mathbf{1}_{[1/n,2/n]}$ satisfies $f_n\to0$ a.e., in measure, almost uniformly, but $\norm{f_n}_p^p=n^{p-1}\nrightarrow0$.
\item Typewriter intervals on $[0,1]$: if $f_n=\mathbf{1}_{I_n}$ and $\lambda(I_n)\to0$ rowwise, then $\norm{f_n}_p\to0$ and $f_n\to0$ in measure, but $(f_n(x))$ converges nowhere.
\item On $(\R,\lambda)$, $f_n=\mathbf{1}_{[n,n+1]}$ satisfies $f_n\to0$ pointwise but $\lambda\{\abs{f_n}\geq1/2\}=1$ for all $n$, so $f_n\nrightarrow0$ in measure.
\item On $(\R,\lambda)$, $f_n=\mathbf{1}_{[n,n+1]}$ is not almost uniformly convergent to $0$: if $\lambda(E)<1$ and $f_n\rightrightarrows0$ on $\R\setminus E$, then $[n,n+1]\subseteq E$ for all large $n$, impossible.
\end{enumerate}
\end{method}

\begin{proposition} \label{prop:finite-measure-integral-metric-for-measure-convergence} \hlt{An Integral Metric for Convergence in Measure on Finite Measure Spaces}
Let $\mu(X)<+\infty$ and $r(h)\defeq\int_X\abs{h}(1+\abs{h})^{-1}\,d\mu$. Then $f_n\to f$ in measure iff $r(f_n-f)\to0$.
\begin{proof}
\begin{enumerate}[label=\roman*.]
\setlength{\itemsep}{0pt}
\item[$(\Rightarrow)$] Set $h_n\defeq f_n-f$. For $\eta>0$ and $A_n\defeq\{\abs{h_n}\geq\eta\}$, $r(h_n)\leq\mu(A_n)+\eta\mu(X)$.\\
Let $n\to\infty$, then $\eta\downarrow0$.
\item[$(\Leftarrow)$] For $\alpha>0$, $\alpha(1+\alpha)^{-1}\mathbf{1}_{\{\abs{h_n}\geq\alpha\}}\leq\abs{h_n}(1+\abs{h_n})^{-1}$, hence $\mu\{\abs{h_n}\geq\alpha\}\leq(1+\alpha)\alpha^{-1}r(h_n)\to0$.
\end{enumerate}
\hlt{Structural Commentary}: The bounded error transform metrizes convergence in measure on finite measure spaces.\\
\hlt{Key Insight}: Small integral of the bounded error transform is equivalent to small large-error sets.
\end{proof}
\end{proposition}

\begin{corollary} \label{cor:fatou-and-dct-with-convergence-in-measure} \hlt{Fatou and Dominated Convergence with In-Measure Hypotheses}
\begin{enumerate}[label=(\roman*)]
\setlength{\itemsep}{0pt}
\item $f_n\geq0$ and $f_n\to f$ in measure $\Rightarrow \int_X f\,d\mu\leq\liminf_n\int_Xf_n\,d\mu$.
\item $f_n\to f$ in measure and $\abs{f_n}\leq g\in\mathscr{L}^1(\mu)$ $\Rightarrow \int_Xf_n\,d\mu\to\int_Xf\,d\mu$.
\end{enumerate}
\begin{proof}
\begin{enumerate}[label=(\roman*)]
\setlength{\itemsep}{0pt}
\item Choose $f_{n_k}$ with $\int f_{n_k}\,d\mu\to\liminf_n\int f_n\,d\mu$. By Theorem \ref{thm:riesz-subsequence-cauchy-measure} and uniqueness, some $f_{n_{k_j}}\to f$ a.e.; Fatou gives $\int f\,d\mu\leq\liminf_j\int f_{n_{k_j}}\,d\mu=\liminf_n\int f_n\,d\mu$.
\item Theorem \ref{thm:dominated-measure-to-lp} with $p=1$ gives $\norm{f_n-f}_1\to0$, hence $\abs{\int_X(f_n-f)\,d\mu}\leq\norm{f_n-f}_1\to0$.
\end{enumerate}
\hlt{Structural Commentary}: In-measure hypotheses reduce to a.e. subsequences or $L^1$ convergence.\\
\hlt{Key Insight}: Fatou and DCT survive convergence in measure through subsequences and domination.
\end{proof}
\end{corollary}

\begin{theorem} \label{thm:vitali-convergence-lp} \hlt{Vitali Convergence Criterion for $L^p$}\\
Let $1\leq p<\infty$, $f_n\in L^p(\mu)$, and let $f$ be measurable. The following are equivalent:
\begin{enumerate}[label=(\roman*)]
\setlength{\itemsep}{0pt}
\item $f\in L^p(\mu)$ and $\norm{f_n-f}_p\to0$.
\item The following hold:
\begin{enumerate}[label=(\alph*)]
\setlength{\itemsep}{0pt}
\item $f_n\to f$ in measure.
\item $\forall\varepsilon>0$ $\exists E_\varepsilon\in\mathcal{M}$ with $\mu(E_\varepsilon)<+\infty$ such that $F\cap E_\varepsilon=\emptyset\Rightarrow\sup_n\int_F\abs{f_n}^p\,d\mu<\varepsilon^p$.
\item $\forall\varepsilon>0$ $\exists\delta(\varepsilon)>0$ such that $\mu(E)<\delta(\varepsilon)\Rightarrow\sup_n\int_E\abs{f_n}^p\,d\mu<\varepsilon^p$.
\end{enumerate}
\end{enumerate}
\begin{proof}
\begin{enumerate}[label=(\roman*)]
\setlength{\itemsep}{0pt}
\item Assume (i). Proposition \ref{prop:lp-convergence-implies-in-measure} gives (ii)(a). Also $(f_n)$ is $L^p$-Cauchy.\\
By Proposition \ref{prop:lp-cauchy-local-mass-control}, for each $\varepsilon>0$ there are $E_\varepsilon\in\mathcal M$ with $\mu(E_\varepsilon)<+\infty$ and $\delta(\varepsilon)>0$ such that $F\cap E_\varepsilon=\emptyset\Rightarrow\sup_n\norm{f_n\mathbf{1}_F}_p<\varepsilon$ and $\mu(E)<\delta(\varepsilon)\Rightarrow\sup_n\norm{f_n\mathbf{1}_E}_p<\varepsilon$; raising to $p$ gives (ii)(b),(ii)(c).
\item Assume (ii). Fix $\eta>0$, choose $E_\eta$ from (ii)(b) and $\delta(\eta)$ from (ii)(c), and put $F=X\setminus E_\eta$.\\
Then $\norm{(f_n-f_m)\mathbf{1}_F}_p\leq\norm{f_n\mathbf{1}_F}_p+\norm{f_m\mathbf{1}_F}_p<2\eta$.\\
If $\mu(E_\eta)>0$, set $\alpha\defeq\eta\mu(E_\eta)^{-1/p}$ and $H_{nm}\defeq E_\eta\cap\{\abs{f_n-f_m}\geq\alpha\}$.\\
By (ii)(a), $(f_n)$ is Cauchy in measure, so $\mu(H_{nm})<\delta(\eta)$ for large $m,n$. On $E_\eta\setminus H_{nm}$, $\abs{f_n-f_m}<\alpha$; on $H_{nm}$, use $\abs{f_n-f_m}\leq\abs{f_n}+\abs{f_m}$. Hence $\norm{(f_n-f_m)\mathbf{1}_{E_\eta}}_p\leq\alpha\mu(E_\eta)^{1/p}+\norm{f_n\mathbf{1}_{H_{nm}}}_p+\norm{f_m\mathbf{1}_{H_{nm}}}_p<3\eta$ by (ii)(c). If $\mu(E_\eta)=0$, the $E_\eta$ term is $0$. Thus $\norm{f_n-f_m}_p<5\eta$ eventually, so $(f_n)$ is $L^p$-Cauchy.\\
Riesz--Fischer gives $f_n\to h$ in $L^p$; Proposition \ref{prop:lp-convergence-implies-in-measure} and Corollary \ref{cor:cauchy-in-measure-completion-uniqueness} give $h=f$ a.e., so (i) holds.
\end{enumerate}
\hlt{Structural Commentary}: Vitali = in-measure convergence + tightness + uniform small-set control.\\
\hlt{Key Insight}: Tightness and small-set control convert in-measure convergence into an $L^p$ Cauchy estimate.
\end{proof}
\end{theorem}

\begin{corollary} \label{cor:domination-implies-vitali-side-conditions} \hlt{Domination Gives Tail and Small-Set Control}\\
Let $1\leq p<\infty$, $g\in L^p(\mu)$, and $\abs{f_n}\leq g$ a.e. for all $n$. Then, for every $\varepsilon>0$:
\begin{enumerate}[label=(\roman*)]
\setlength{\itemsep}{0pt}
\item $\exists E_\varepsilon\in\mathcal M$ with $\mu(E_\varepsilon)<+\infty$ such that $F\cap E_\varepsilon=\emptyset\Rightarrow\sup_n\int_F\abs{f_n}^p\,d\mu<\varepsilon^p$.
\item $\exists\delta(\varepsilon)>0$ such that $\mu(E)<\delta(\varepsilon)\Rightarrow\sup_n\int_E\abs{f_n}^p\,d\mu<\varepsilon^p$.
\end{enumerate}
\begin{proof}
\begin{enumerate}[label=(\roman*)]
\setlength{\itemsep}{0pt}
\item Proposition \ref{prop:simple-density-support-and-tails-lp} gives $E_\varepsilon$ with $\mu(E_\varepsilon)<+\infty$ and $F\cap E_\varepsilon=\emptyset\Rightarrow\int_F\abs g^p\,d\mu<\varepsilon^p$.\\
Since $\abs{f_n}^p\leq\abs g^p$ a.e., $\sup_n\int_F\abs{f_n}^p\,d\mu<\varepsilon^p$.
\item Lemma \ref{lem:lp-integral-absolute-continuity} gives $\delta(\varepsilon)>0$ with $\mu(E)<\delta(\varepsilon)\Rightarrow\int_E\abs g^p\,d\mu<\varepsilon^p$.\\
Since $\abs{f_n}^p\leq\abs g^p$ a.e., $\sup_n\int_E\abs{f_n}^p\,d\mu<\varepsilon^p$.
\end{enumerate}
\hlt{Structural Commentary}: One $L^p$ envelope supplies both tail tightness and small-set control.\\
\hlt{Key Insight}: Domination transfers both estimates from $g$ to every $f_n$.
\end{proof}
\end{corollary}

\begin{proposition} \label{prop:composition-preserves-pointwise-and-almost-uniform-modes} \hlt{Composition and Pointwise Modes}\\
Let $\varphi:\R\to\R$.
\begin{enumerate}[label=(\roman*)]
\setlength{\itemsep}{0pt}
\item $\varphi$ continuous, $f_n\to f$ a.e. $\Rightarrow \varphi\circ f_n\to\varphi\circ f$ a.e.
\item $\varphi$ uniformly continuous preserves uniform, a.u., and in-measure convergence.
\item If $\varphi$ is discontinuous, composition need not preserve a.e. convergence.
\item Non-uniform continuity can destroy uniform, even in-measure, convergence.
\end{enumerate}
\begin{proof}
\begin{enumerate}[label=(\roman*)]
\setlength{\itemsep}{0pt}
\item Outside a null set, $f_n(x)\to f(x)$; continuity gives $\varphi(f_n(x))\to\varphi(f(x))$.
\item Given $\varepsilon>0$, choose $\delta>0$ with $\abs{s-t}<\delta\Rightarrow\abs{\varphi(s)-\varphi(t)}<\varepsilon$. Uniform and almost-uniform convergence use $\sup\abs{f_n-f}<\delta$ on $X$ or $X\setminus E_\eta$; in measure, $\{\abs{\varphi(f_n)-\varphi(f)}\geq\varepsilon\}\subseteq\{\abs{f_n-f}\geq\delta\}$.
\item If $\varphi$ is discontinuous at $a$, choose $t_n\to a$ and $\alpha>0$ with $\abs{\varphi(t_n)-\varphi(a)}\geq\alpha$.\\
On $X=\{\ast\}$, $f_n(\ast)=t_n\to a=f(\ast)$, but $\varphi(f_n(\ast))\nrightarrow\varphi(f(\ast))$.
\item If $\varphi$ is not uniformly continuous, choose $a_n,b_n$ with $\abs{a_n-b_n}\to0$ and $\abs{\varphi(a_n)-\varphi(b_n)}\geq\alpha>0$.\\
On $(\N,\#)$, set $f(k)=b_k$ and $f_n(k)=a_n$ for $k=n$, $f_n(k)=b_k$ for $k\neq n$.\\
Then $\norm{f_n-f}_\infty\to0$, but $\#\{\abs{\varphi(f_n)-\varphi(f)}\geq\alpha\}=1$ for all $n$.
\end{enumerate}
\hlt{Structural Commentary}: Continuity is pointwise; uniform continuity gives one global $\delta$.\\
\hlt{Key Insight}: Composition needs the same uniformity as the convergence mode.
\end{proof}
\end{proposition}

\begin{proposition} \label{prop:composition-preserves-convergence-in-measure} \hlt{Composition and Convergence in Measure}\\
Let $f_n,f:X\to\R$ be measurable with $f_n\to f$ in measure, and let $\varphi:\R\to\R$.
\begin{enumerate}[label=(\roman*)]
\setlength{\itemsep}{0pt}
\item If $\varphi$ is uniformly continuous, then $\mu\{\abs{\varphi(f_n)-\varphi(f)}\geq\varepsilon\}\to0$ for every $\varepsilon>0$.
\item If $\mu(X)<+\infty$ and $\varphi$ is continuous, then $\mu\{\abs{\varphi(f_n)-\varphi(f)}\geq\varepsilon\}\to0$ for every $\varepsilon>0$.
\end{enumerate}
\begin{proof}
\begin{enumerate}[label=(\roman*)]
\setlength{\itemsep}{0pt}
\item Fix $\varepsilon>0$ and choose $\delta>0$ with $\abs{s-t}<\delta\Rightarrow\abs{\varphi(s)-\varphi(t)}<\varepsilon$.\\
Then $\{\abs{\varphi(f_n)-\varphi(f)}\geq\varepsilon\}\subseteq\{\abs{f_n-f}\geq\delta\}$, so the measures tend to $0$.
\item Fix $\varepsilon,\eta>0$. Since $\mu(X)<+\infty$ and $f$ is finite-valued, choose $M$ with $\mu\{\abs f>M\}<\eta$. Uniform continuity of $\varphi$ on $[-M-1,M+1]$ gives $0<\delta<1$ such that $s,t\in[-M-1,M+1]$ and $\abs{s-t}<\delta\Rightarrow\abs{\varphi(s)-\varphi(t)}<\varepsilon$. Hence $\{\abs{\varphi(f_n)-\varphi(f)}\geq\varepsilon\}\subseteq\{\abs f>M\}\cup\{\abs{f_n-f}\geq\delta\}$, so $\limsup_n\mu\{\abs{\varphi(f_n)-\varphi(f)}\geq\varepsilon\}\leq\eta\downarrow0$.
\end{enumerate}
\hlt{Structural Commentary}: Use global uniformity, or truncate to a bounded range.\\
\hlt{Key Insight}: On finite measure spaces, measurable finite-valued functions are almost bounded.
\end{proof}
\end{proposition}

\begin{example} \label{ex:continuous-composition-in-measure-needs-hypotheses} \hlt{Continuity Alone Is Not Enough on Infinite Measure}\\
On $(\R,\mathcal{B}(\R),\lambda)$, set $f(x)=x$, $f_n(x)=x+1/n$, and $\varphi(t)=t^2$.\\
Then $\norm{f_n-f}_\infty=1/n\to0$, hence $f_n\rightrightarrows f$ and $f_n\to f$ in measure.\\
However $\varphi(f_n)-\varphi(f)=2x/n+1/n^2$; for $E_n\defeq\{\abs{\varphi(f_n)-\varphi(f)}\geq1\}$, $[n/2,+\infty)\subseteq E_n$, so $\lambda(E_n)=+\infty$.\\
Thus $\varphi\circ f_n\nrightarrow\varphi\circ f$ in measure.
\end{example}

\begin{proposition} \label{prop:lp-composition-growth-and-continuity} \hlt{Growth Conditions for $L^p$ Composition}\\
Let $1\leq p<\infty$, let $\varphi:\R\to\R$ be measurable, and let $f_n,f:X\to\R$ be measurable.
\begin{enumerate}[label=(\roman*)]
\setlength{\itemsep}{0pt}
\item If $\abs{\varphi(t)}\leq A\abs t+B$ for all $t$ and some $A,B\geq0$, $\mu(X)<+\infty$, and $f\in L^p(\mu)$, then $\varphi\circ f\in L^p(\mu)$.
\item If $\abs{\varphi(t)}\leq A\abs t$ for all $t$ and some $A\geq0$, and $f\in L^p(\mu)$, then $\varphi\circ f\in L^p(\mu)$.
\item If $\mu(X)<+\infty$, $\varphi$ is uniformly continuous, $f_n,f\in L^p(\mu)$, and $\norm{f_n-f}_p\to0$, then $\norm{\varphi\circ f_n-\varphi\circ f}_p\to0$.
\item If $\varphi$ is Lipschitz, $\varphi(0)=0$, $f_n,f\in L^p(\mu)$, and $\norm{f_n-f}_p\to0$, then $\norm{\varphi\circ f_n-\varphi\circ f}_p\to0$.
\end{enumerate}
\begin{proof}
\begin{enumerate}[label=(\roman*)]
\setlength{\itemsep}{0pt}
\item $\abs{\varphi(f)}^p\leq2^{p-1}(A^p\abs f^p+B^p)$, so $\norm{\varphi\circ f}_p^p\leq2^{p-1}(A^p\norm f_p^p+B^p\mu(X))<+\infty$.
\item $\abs{\varphi(f)}^p\leq A^p\abs f^p$, so $\norm{\varphi\circ f}_p\leq A\norm f_p<+\infty$.
\item Uniform continuity gives constants $C,D\geq0$ with $\abs{\varphi(t)}\leq C\abs t+D$ and $\abs{\varphi(s)-\varphi(t)}\leq C\abs{s-t}+D$.\\
Hence all compositions lie in $L^p$ by (i). Let $h_n\defeq\abs{f_n-f}$.\\
Given $\eta>0$, choose $\delta>0$ with $\abs{s-t}<\delta\Rightarrow\abs{\varphi(s)-\varphi(t)}<\eta$ and set $H_n\defeq\{h_n\geq\delta\}$.\\
Then $\int_X\abs{\varphi(f_n)-\varphi(f)}^p\,d\mu\leq\eta^p\mu(X)+2^{p-1}\int_{H_n}(C^ph_n^p+D^p)\,d\mu$.\\
Since $\int h_n^p\,d\mu\to0$ and $\mu(H_n)\leq\delta^{-p}\int h_n^p\,d\mu\to0$, take $\limsup_n$ and then $\eta\downarrow0$.
\item If $L$ is Lipschitz, $\abs{\varphi(f)}\leq L\abs f$ and $\norm{\varphi\circ f_n-\varphi\circ f}_p\leq L\norm{f_n-f}_p\to0$.
\end{enumerate}
\hlt{Structural Commentary}: Growth gives membership; continuity estimates transfer $L^p$ convergence.\\
\hlt{Key Insight}: Finite measure makes small uniform errors integrable; Lipschitz bounds give norm control.
\end{proof}
\end{proposition}

\begin{example} \label{ex:uniform-continuity-composition-lp-needs-finite-measure} \hlt{Uniform Continuity Does Not Suffice on Infinite Measure}\\
On $(\R,\lambda)$, fix $1\leq p<\infty$, set $f=0$, $f_n\defeq n^{-1}\mathbf{1}_{[0,n^{p/2}]}$, and $\varphi(t)\defeq\sqrt{\abs t}$.\\
Then $\varphi$ is uniformly continuous, $\varphi(0)=0$, $\norm{f_n-f}_p^p=n^{-p/2}\to0$, but $\norm{\varphi\circ f_n-\varphi\circ f}_p^p=1$ for all $n$.\\
Thus uniform continuity need not preserve $L^p$ convergence on infinite measure spaces.
\end{example}

\begin{lemma} \label{lem:lp-convergence-continuous-composition-in-measure} \hlt{$L^p$ Convergence Allows Continuous Composition in Measure}\\
Let $1\leq p<\infty$, $f_n,f\in L^p(\mu)$, $\norm{f_n-f}_p\to0$, and let $\varphi:\R\to\R$ be continuous.\\
Then $\varphi\circ f_n\to\varphi\circ f$ in measure.
\begin{proof}
Fix $\alpha,\eta>0$. Choose $M$ with $\mu\{\abs f>M\}<\eta$. By uniform continuity of $\varphi$ on $[-M-1,M+1]$, choose $0<\delta<1$ such that $s,t\in[-M-1,M+1]$ and $\abs{s-t}<\delta\Rightarrow\abs{\varphi(s)-\varphi(t)}<\alpha$.\\
Then $\{\abs{\varphi(f_n)-\varphi(f)}\geq\alpha\}\subseteq\{\abs f>M\}\cup\{\abs{f_n-f}\geq\delta\}$, so $\limsup_n\mu\{\abs{\varphi(f_n)-\varphi(f)}\geq\alpha\}\leq\eta$. Let $\eta\downarrow0$.\\
\hlt{Structural Commentary}: $L^p$ convergence localizes to a large bounded range where continuity is uniform.\\
\hlt{Key Insight}: Norm convergence gives in-measure control; bounded-range continuity transfers it.
\end{proof}
\end{lemma}

\begin{theorem} \label{thm:composition-growth-criteria-lp} \hlt{Composition Growth Criteria in $L^p$}\\
Let $1\leq p<\infty$ and let $\varphi:\R\to\R$ be continuous.
\begin{enumerate}[label=(\roman*)]
\setlength{\itemsep}{0pt}
\item If $\mu(X)<+\infty$ and $\exists K>0$ such that $\abs{\varphi(t)}\leq K\abs t$ for $\abs t\geq K$, then $f\in L^p(\mu)\Rightarrow\varphi\circ f\in L^p(\mu)$.
\item Under (i), if $f_n,f\in L^p(\mu)$ and $\norm{f_n-f}_p\to0$, then $\norm{\varphi\circ f_n-\varphi\circ f}_p\to0$.
\item If $\exists K\geq0$ such that $\abs{\varphi(t)}\leq K\abs t$ for all $t$, then $f\in L^p(\mu)\Rightarrow\varphi\circ f\in L^p(\mu)$ on any measure space.
\item Under (iii), if $f_n,f\in L^p(\mu)$ and $\norm{f_n-f}_p\to0$, then $\norm{\varphi\circ f_n-\varphi\circ f}_p\to0$ on any measure space.
\end{enumerate}
\begin{proof}
\begin{enumerate}[label=(\roman*)]
\setlength{\itemsep}{0pt}
\item Let $M\defeq\sup_{\abs t\leq K}\abs{\varphi(t)}$ and $C\defeq K+M$.\\
Then $\abs{\varphi(t)}\leq C(1+\abs t)$ for all $t$. hence $\abs{\varphi(f)}^p\leq2^{p-1}C^p(1+\abs f^p)\in L^1(\mu)$.
\item Let $C$ be as in (i). Lemma \ref{lem:lp-convergence-continuous-composition-in-measure} gives $\varphi\circ f_n\to\varphi\circ f$ in measure.\\
Since $(f_n)$ is $L^p$-Cauchy, Proposition \ref{prop:lp-cauchy-local-mass-control} and $\abs{\varphi(f_n)}^p\leq2^{p-1}C^p(1+\abs{f_n}^p)$ give the finite-tail estimate with $E_\varepsilon=X$ and the small-set estimate. Theorem \ref{thm:vitali-convergence-lp} gives $\norm{\varphi\circ f_n-\varphi\circ f}_p\to0$.
\item $\abs{\varphi(f)}^p\leq K^p\abs f^p\in L^1(\mu)$.
\item Lemma \ref{lem:lp-convergence-continuous-composition-in-measure} gives $\varphi\circ f_n\to\varphi\circ f$ in measure. Since $(f_n)$ is $L^p$-Cauchy and $\abs{\varphi(f_n)}^p\leq K^p\abs{f_n}^p$, Proposition \ref{prop:lp-cauchy-local-mass-control} gives the finite-tail and small-set estimates for $(\varphi\circ f_n)$. Theorem \ref{thm:vitali-convergence-lp} gives $\norm{\varphi\circ f_n-\varphi\circ f}_p\to0$.
\end{enumerate}
\hlt{Structural Commentary}: Local growth needs finite measure; global growth does not.\\
\hlt{Key Insight}: Vitali converts in-measure composition plus inherited mass control into $L^p$ convergence.
\end{proof}
\end{theorem}

\begin{method} \label{meth:composition-growth-counterexamples} \hlt{Atomic Counterexamples for Composition Growth}
\begin{enumerate}[label=(\roman*)]
\setlength{\itemsep}{0pt}
\item If finite-measure tail growth fails, choose $\abs{t_n}\geq n$ with $\abs{\varphi(t_n)}>n^2\abs{t_n}$. On $\N$, set $\mu(\{n\})=(n^2\abs{t_n}^p)^{-1}$ and $f(n)=t_n$. Then $\mu(\N)<+\infty$, $\int\abs f^p\,d\mu=\sum_n n^{-2}<+\infty$, but $\int\abs{\varphi(f)}^p\,d\mu\geq\sum_n n^{2p-2}=+\infty$.
\item If the same tail growth fails, choose $\abs{t_n}\geq n$ with $\abs{\varphi(t_n)-\varphi(0)}\geq n\abs{t_n}$. With $\mu(\{n\})=(n^p\abs{t_n}^p)^{-1}$, $f=0$, and $f_n=t_n\mathbf{1}_{\{n\}}$, one has $\mu(\N)<+\infty$, $\norm{f_n-f}_p=1/n\to0$, but $\norm{\varphi\circ f_n-\varphi\circ f}_p\geq1$.
\item If global linear growth fails and $\varphi(0)=0$, the constructions in (i)--(ii), without requiring $\mu(\N)<+\infty$, give membership and convergence counterexamples on arbitrary measure spaces.
\item If $\varphi(0)\neq0$, then on $(\N,\#)$, $0\in L^p$ but $\varphi\circ0$ is a nonzero constant and hence not in $L^p$; the constant sequence $f_n=0$ gives the convergence failure.
\end{enumerate}
\end{method}
